\lecture{1.7}{Regression as Conditional Expectation Estimation}
\label{lec:1-7}

\subsection{Regression as Estimation of $\E(Y\given A,W)$}

In this course, regression will be used primarily as a tool for estimating conditional expectations. This is broader than the common habit of treating regression coefficients as the main scientific output.

\begin{definition}[Regression function]
For outcome $Y$, treatment or exposure $A$, and covariates $W$, the regression function is the conditional mean
\[
\Qbar(a,w)=\E(Y\given A=a,W=w).
\]
An estimator of this function will be denoted $\hat Q_n(a,w)$ or $\Qbar_n(a,w)$.
\end{definition}

The regression function is a machine: input a subgroup specification $(a,w)$, output the mean of $Y$ among units in that subgroup.

\begin{estimationbox}
Regression workflow for estimating $\E(Y\given A,W)$:
\begin{enumerate}
  \item Choose a regression method and a set of predictors involving $A$ and $W$.
  \item Fit the regression using the observed data $(Y_i,A_i,W_i)$.
  \item Obtain an estimated function $\Qbar_n(a,w)$.
  \item Evaluate $\Qbar_n(a,w)$ at covariate and treatment values of scientific interest.
\end{enumerate}
\end{estimationbox}

\subsection{Linear Regression}

A main-terms linear regression model for a scalar covariate $W$ is
\[
\E(Y\given A,W)=\beta_0+\beta_1A+\beta_2W.
\]
After fitting the model, the estimated regression function is
\[
\Qbar_n(a,w)=\hat\beta_0+\hat\beta_1a+\hat\beta_2w.
\]

For example, if the fitted coefficients are
\[
\hat\beta_0=3.27,\qquad
\hat\beta_1=0.17,\qquad
\hat\beta_2=-0.18,
\]
then
\[
\Qbar_n(a,w)=3.27+0.17a-0.18w.
\]
The estimated mean outcome for untreated people in age group 4 is
\[
\Qbar_n(0,4)=3.27+0.17(0)-0.18(4)=2.55.
\]
The estimated mean outcome for treated people in age group 7 is
\[
\Qbar_n(1,7)=3.27+0.17(1)-0.18(7)=2.18.
\]

The coefficient $\hat\beta_1$ is the modeled difference between treatment groups holding $W$ fixed by one unit of the model. Whether that coefficient is itself a meaningful scientific estimand depends on the question and on the model. The safer general view is that the fitted model estimates a conditional mean function.

\subsection{Logistic Regression}

When the outcome is binary, a conditional mean must lie between 0 and 1. Logistic regression ensures this by modeling the log-odds and then transforming back to a probability.

\begin{definition}[Expit function]
The expit function is
\[
\expit(x)=\frac{\exp(x)}{1+\exp(x)}
=\frac{1}{1+\exp(-x)}.
\]
It maps every real number to a value in $(0,1)$.
\end{definition}

A main-terms logistic regression model is
\[
\E(Y\given A,W)
=\Prob(Y=1\given A,W)
=\expit(\gamma_0+\gamma_1A+\gamma_2W).
\]
After fitting,
\[
\Qbar_n(a,w)=\expit(\hat\gamma_0+\hat\gamma_1a+\hat\gamma_2w).
\]

If logistic regression is used to estimate the treatment mechanism rather than the outcome regression, one might write
\[
g(w)=\Prob(A=1\given W=w)
=\expit(\alpha_0+\alpha_1w),
\]
with fitted estimator
\[
g_n(w)=\expit(\hat\alpha_0+\hat\alpha_1w).
\]
For example, if $\hat\alpha_0=-0.44$ and $\hat\alpha_1=0.07$, then
\[
g_n(2)=\expit(-0.44+0.07\cdot 2)\approx 0.43,
\]
and
\[
g_n(5)=\expit(-0.44+0.07\cdot 5)\approx 0.48.
\]

\subsection{Generalized Linear Models}

Linear and logistic regression are examples of generalized linear models (GLMs). A GLM specifies a conditional mean through a link function:
\[
h\{\E(Y\given X)\}=\eta(X),
\]
where $X$ denotes predictors, $h$ is the link function, and $\eta(X)$ is a linear predictor.

For linear regression, $h(u)=u$, so
\[
\E(Y\given X)=\eta(X).
\]
For logistic regression, $h(u)=\logit(u)$, so
\[
\logit\{\E(Y\given X)\}=\eta(X),
\qquad
\E(Y\given X)=\expit\{\eta(X)\}.
\]

In causal inference applications, GLMs are often used as nuisance regressions: intermediate models needed to estimate a target causal parameter. The final target is often not a regression coefficient.

\subsection{Main-Terms Models}

A main-terms model includes predictors additively on the scale of the linear predictor. With scalar $W$,
\[
\eta(A,W)=\beta_0+\beta_1A+\beta_2W.
\]
For several covariates $W=(W_1,\ldots,W_p)$,
\[
\eta(A,W)=\beta_0+\beta_1A+\beta_2W_1+\cdots+\beta_{p+1}W_p.
\]

Main-terms models borrow information across subgroups by imposing a simple structure. In the linear model above, the difference between $A=1$ and $A=0$ is the same at every value of $W$:
\[
\Qbar_n(1,w)-\Qbar_n(0,w)=\hat\beta_1.
\]
That restriction can be useful when data are limited, but it can be wrong if treatment-outcome associations vary across covariate values.

\subsection{Interaction Models}

Interaction models allow the relationship between one predictor and the outcome to depend on another predictor. A linear model with an $A$ by $W$ interaction is
\[
\E(Y\given A,W)=\beta_0+\beta_1A+\beta_2W+\beta_3AW.
\]
The fitted function is
\[
\Qbar_n(a,w)=\hat\beta_0+\hat\beta_1a+\hat\beta_2w+\hat\beta_3aw.
\]
Now the fitted difference between treatment groups at covariate value $w$ is
\[
\Qbar_n(1,w)-\Qbar_n(0,w)=\hat\beta_1+\hat\beta_3w.
\]
Thus the treatment contrast on the outcome scale is allowed to vary with $W$.

Polynomial terms are another way to add flexibility:
\[
\eta(A,W)=\beta_0+\beta_1A+\beta_2W+\beta_3W^2.
\]
The choice of regression formula determines what patterns can and cannot be represented by the fitted conditional mean.

\subsection{Using \texttt{glm} in R}

The R function \texttt{glm} fits generalized linear models. Suppose a data frame \texttt{covid19} contains:
\begin{itemize}
  \item \texttt{out}: an outcome variable $Y$;
  \item \texttt{treat}: a binary treatment variable $A$;
  \item \texttt{age\_grp}: an age group covariate $W$.
\end{itemize}

A main-terms linear regression for $\E(Y\given A,W)$ is fit by:
\begin{lstlisting}[language=R]
fit1 <- glm(out ~ treat + age_grp, data = covid19)
coef(fit1)
\end{lstlisting}

If the coefficient output is approximately
\begin{lstlisting}
(Intercept)      treat    age_grp
     3.2714     0.1692    -0.1780
\end{lstlisting}
then the fitted conditional mean is
\[
\Qbar_n(a,w)=3.2714+0.1692a-0.1780w.
\]

A main-terms logistic regression for the treatment probability $\Prob(A=1\given W)$ is fit by:
\begin{lstlisting}[language=R]
fit2 <- glm(treat ~ age_grp, data = covid19, family = binomial())
coef(fit2)
\end{lstlisting}

If the fitted coefficients are approximately $-0.4391$ for the intercept and $0.0729$ for \texttt{age\_grp}, then
\[
g_n(w)=\expit(-0.4391+0.0729w).
\]

\subsection{Prediction with \texttt{newdata}}

After fitting a regression, \texttt{predict} evaluates the fitted conditional mean at specified covariate values. For example:
\begin{lstlisting}[language=R]
aw_values <- data.frame(
  treat = c(0, 1),
  age_grp = c(4, 7)
)

predict(fit1, newdata = aw_values, type = "response")
\end{lstlisting}

The object \texttt{aw\_values} contains two subgroup specifications: $(A=0,W=4)$ and $(A=1,W=7)$. The call to \texttt{predict} returns the fitted means $\Qbar_n(0,4)$ and $\Qbar_n(1,7)$.

For logistic regression, \texttt{type = "response"} returns probabilities rather than log-odds:
\begin{lstlisting}[language=R]
w_values <- data.frame(age_grp = c(2, 5))
predict(fit2, newdata = w_values, type = "response")
\end{lstlisting}

These are estimates of $\Prob(A=1\given W=2)$ and $\Prob(A=1\given W=5)$ under the fitted logistic model.

\subsection{Regression Output as Subgroup Mean Estimates}

Regression output should be translated back into conditional mean estimates. Coefficients are ingredients for constructing the function $\Qbar_n$ or $g_n$.

\begin{keyidea}
A fitted regression is an estimated function. It takes subgroup-defining variables as input and returns an estimated subgroup mean or subgroup probability.
\end{keyidea}

Sometimes we need a function such as
\[
w\mapsto \E(Y\given A=1,W=w)
\]
rather than the full function $(a,w)\mapsto \E(Y\given A=a,W=w)$. There are two common approaches.

\paragraph{Pooling.}
Fit a regression using all observations and include $A$ and $W$ as predictors. If
\[
\Qbar_n(a,w)=\hat\beta_0+\hat\beta_1a+\hat\beta_2w,
\]
then estimate the treated conditional mean by setting $a=1$:
\[
\Qbar_{n,1}(w)=\Qbar_n(1,w)=\hat\beta_0+\hat\beta_1+\hat\beta_2w.
\]

\paragraph{Stratifying.}
Restrict the data to observations with $A=1$ and regress $Y$ on $W$ in that subset. A main-terms stratified fit might yield
\[
\Qbar_{n,1}(w)=\hat\eta_0+\hat\eta_1w.
\]

Pooling borrows information across treatment groups through the model structure. Stratifying avoids imposing a common structure across treatment groups but may use less data. Later causal estimators will use both ideas, often with more flexible regression methods.

\begin{remark}
Regression coefficients can be scientifically meaningful in special settings, but the general causal inference workflow is different: define the target estimand first, identify it in terms of observed-data quantities, and then use regression as needed to estimate those quantities.
\end{remark}
